\chapter{4选1数据选择器（MUX）}

\section{第一步：解决什么问题}

想象一个场景：你需要从4个数据源中选择1个传给后续电路。

\begin{itemize}
  \item 不是同时传4个——那需要4条数据通路
  \item 而是：用2根选择线指定"要哪一个"，1根数据线输出
\end{itemize}

\begin{keypoint}[MUX的本质]
数据选择器（Multiplexer, MUX）= 一个受控的多路开关

$2^n$ 个数据输入 $\to$ $n$ 条选择线 $\to$ 1个输出

选择线决定"哪一路数据通向输出"
\end{keypoint}

\section{第二步：输入输出关系}

4选1 MUX：

\begin{itemize}
  \item \textbf{数据输入}：$D_0, D_1, D_2, D_3$
  \item \textbf{选择输入}：$S_1, S_0$（2位，选择哪个数据输入）
  \item \textbf{输出}：$Y$
  \item \textbf{功能}：$Y = D_{(S_1 S_0)}$
\end{itemize}

\section{第三步：真值表}

\begin{center}
\begin{tabular}{|c|c||c|}
\hline
$S_1$ & $S_0$ & $Y$ \\
\hline
0 & 0 & $D_0$ \\
0 & 1 & $D_1$ \\
1 & 0 & $D_2$ \\
1 & 1 & $D_3$ \\
\hline
\end{tabular}
\end{center}

\section{第四步：逻辑推导}

从真值表可以写出：

\begin{equation}
  Y = \overline{S_1} \cdot \overline{S_0} \cdot D_0
    + \overline{S_1} \cdot S_0 \cdot D_1
    + S_1 \cdot \overline{S_0} \cdot D_2
    + S_1 \cdot S_0 \cdot D_3
\end{equation}

\textbf{关键观察}：MUX的每个选择组合对应一个\textbf{乘积项}。

\begin{itemize}
  \item 选择线产生最小项（$\overline{S_1}\cdot\overline{S_0}$, $\overline{S_1}\cdot S_0$等）
  \item 每个数据输入与其对应的最小项AND
  \item 所有结果OR起来
\end{itemize}

\section{第五步：结构化设计}

4选1 MUX内部结构 = 4个3输入AND门 + 1个4输入OR门 + 2个反相器：

\begin{center}
\begin{tikzpicture}[scale=0.85]
  % Inputs
  \node (D0) at (-4,3) {$D_0$};
  \node (D1) at (-4,1) {$D_1$};
  \node (D2) at (-4,-1) {$D_2$};
  \node (D3) at (-4,-3) {$D_3$};
  \node (S1) at (-4,-5) {$S_1$};
  \node (S0) at (-4,-6) {$S_0$};

  % Internal structure is complex, using simplified representation
  \node[draw, minimum width=4cm, minimum height=5cm,
        label=above:MUX 4选1] (mux) at (0,0) {};

  \node (Y) at (3,0) {$Y$};
  \draw[->] (mux.east) -- (Y);

  % Selection notation
  \node at (0,-3.5) {\small $S_1S_0$};
\end{tikzpicture}
\end{center}

\section{第六步：为什么这样设计}

\begin{enumerate}
  \item \textbf{为什么MUX是"万能"的组合逻辑？}

  因为MUX输出 = $\sum$（选择最小项 $\cdot$ 数据输入）。

  如果把数据输入当作"该最小项是否参与"的控制：
  \begin{itemize}
    \item $D_k = 1$：该最小项参与（函数包含此最小项）
    \item $D_k = 0$：该最小项不参与
  \end{itemize}

  所以：\textbf{任何$n$变量逻辑函数都可以用一个$2^n$选1 MUX实现}！

  只需将函数真值表的输出列接到MUX的数据输入端。

  \item \textbf{为什么MUX能节省资源？}

  将"多路数据"变成"一路分时复用"——不需要为每路数据各建一条独立通路。这是时分复用的基本思想。

  \item \textbf{MUX与译码器的关系？}

  译码器产生所有$2^n$个最小项（$n$输入$\to$$2^n$输出），MUX选择其中一个输出。实际上，MUX = 译码器 + AND-OR结构。
\end{enumerate}

\section{第七步：考试常见变式}

\begin{enumerate}
  \item \textbf{变式1：用MUX实现任意逻辑函数（必考！）}

  \textbf{方法}：将函数的变量接选择线，真值表的输出值接数据输入。

  例：实现 $F(A,B,C) = \sum m(0,1,4,6,7)$

  用8选1 MUX：$D_0=1, D_1=1, D_2=0, D_3=0, D_4=1, D_5=0, D_6=1, D_7=1$

  \textbf{降维技巧}：用$2^{n-1}$选1 MUX实现$n$变量函数

  选择线只接$n-1$个变量，第$n$个变量（或其反变量）接到某些数据输入端。

  例：用4选1 MUX实现3变量函数 $F(A,B,C)$

  将A,B接选择线 $S_1,S_0$，C（或其反）接到某些数据输入端。

  \item \textbf{变式2：双4选1 MUX（8选1功能）}

  两片4选1 MUX + 一片2选1 MUX = 8选1 MUX。

  \item \textbf{变式3：带使能端的MUX}

  使能端E=0时，输出固定为0（或高阻态），无论选择线是什么。

  \item \textbf{变式4：MUX扩展}

  $n$选1 MUX可以通过级联构建更大的MUX。例如：4片4选1 MUX + 1片4选1 MUX = 16选1 MUX。
\end{enumerate}

\section{第八步：VHDL实现}

\begin{lstlisting}[caption={4选1 MUX（多种实现方式）}]
-- 方式1: with-select (最简洁)
architecture with_select of mux4to1 is
begin
  with S select
    Y <= D(0) when "00",
         D(1) when "01",
         D(2) when "10",
         D(3) when "11",
         '0'   when others;
end with_select;

-- 方式2: when-else (条件信号赋值)
architecture when_else of mux4to1 is
begin
  Y <= D(0) when S = "00" else
       D(1) when S = "01" else
       D(2) when S = "10" else
       D(3) when S = "11" else
       '0';
end when_else;

-- 方式3: process + case
architecture proc_case of mux4to1 is
begin
  process(S, D)
  begin
    case S is
      when "00" => Y <= D(0);
      when "01" => Y <= D(1);
      when "10" => Y <= D(2);
      when "11" => Y <= D(3);
      when others => Y <= '0';
    end case;
  end process;
end proc_case;
\end{lstlisting}

\begin{lstlisting}[caption={完整4选1 MUX实体}]
library ieee;
use ieee.std_logic_1164.all;

entity mux4to1 is
  port (
    D : in  std_logic_vector(3 downto 0);
    S : in  std_logic_vector(1 downto 0);
    Y : out std_logic
  );
end mux4to1;
\end{lstlisting}

\section{第九步：Testbench验证}

\begin{lstlisting}[caption={4选1 MUX Testbench}]
library ieee;
use ieee.std_logic_1164.all;

entity mux4to1_tb is
end mux4to1_tb;

architecture test of mux4to1_tb is
  signal D : std_logic_vector(3 downto 0);
  signal S : std_logic_vector(1 downto 0);
  signal Y : std_logic;

  component mux4to1 is
    port (
      D : in  std_logic_vector(3 downto 0);
      S : in  std_logic_vector(1 downto 0);
      Y : out std_logic
    );
  end component;
begin
  uut: mux4to1 port map (D, S, Y);

  stimulus: process
  begin
    D <= "1010";  -- 固定数据输入
    S <= "00"; wait for 10 ns;  -- Y = D0 = 0
    S <= "01"; wait for 10 ns;  -- Y = D1 = 1
    S <= "10"; wait for 10 ns;  -- Y = D2 = 0
    S <= "11"; wait for 10 ns;  -- Y = D3 = 1
    wait;
  end process;
end test;
\end{lstlisting}

\section{第十步：如何快速解题}

\begin{examtip}[MUX快速解题法]
\begin{enumerate}
  \item \textbf{用MUX实现逻辑函数}：真值表的输出列 = MUX的数据输入列（这是最快的做法）
  \item \textbf{降维法}：$n$变量函数 $\to$ 用$2^{n-1}$选1 MUX。剩余的1个变量（或其反）接数据输入
  \item \textbf{级联扩展}：小MUX + 小MUX = 大MUX
  \item \textbf{VHDL}：with-select 是最简洁的MUX写法，综合结果就是MUX硬件
\end{enumerate}
\end{examtip}

\textbf{秒杀必考题：用MUX实现逻辑函数}

\noindent 题：用8选1 MUX实现 $F(A,B,C) = \sum m(0, 3, 5, 6)$

\noindent 解：$D_0=1, D_1=0, D_2=0, D_3=1, D_4=0, D_5=1, D_6=1, D_7=0$

\section{变式详解}
\chapter{4选1 MUX——4大变式详解}

\section{变式1：用MUX实现任意逻辑函数（必考！）}

\begin{detailbox}[完整教学]
\textbf{【核心原理】}
MUX的通用表达式：$Y=\overline{S_1}\cdot\overline{S_0}\cdot D_0+\overline{S_1}\cdot S_0\cdot D_1+S_1\cdot\overline{S_0}\cdot D_2+S_1\cdot S_0\cdot D_3$

如果把$S_1,S_0$接函数变量，$D_k$接"该最小项是否参与=1"：$D_k$接1→该最小项参与函数。$D_k$接0→不参与。

\textbf{【结论】任何n变量逻辑函数=一个$2^n$选1 MUX}

\textbf{【例题1】}用8选1 MUX实现$F(A,B,C)=\sum m(0,3,5,6)$

\textbf{解：}$D_0=1,D_1=0,D_2=0,D_3=1,D_4=0,D_5=1,D_6=1,D_7=0$

将m0-m7对应的D分别接1或0。选择线S2,S1,S0接A,B,C。

\textbf{【例题2：降维法】}用4选1 MUX实现3变量函数$F(A,B,C)=\sum m(0,1,4,6,7)$

$2^{3-1}=4$选1 MUX, 变量A,B接S1,S0, C用于产生$D_k$。

\begin{center}
\begin{tabular}{|c|c|c|c|c|}
\hline
AB & m(C=0) & m(C=1) & D输入 & 逻辑 \\
\hline
00 & m0=1 & m1=1 & D0 & 1(与C无关) \\
01 & m2=0 & m3=0 & D1 & 0 \\
10 & m4=1 & m5=0 & D2 & $\overline{C}$ \\
11 & m6=1 & m7=1 & D3 & 1 \\
\hline
\end{tabular}
\end{center}

D0=1, D1=0, D2=$\overline{C}$, D3=1.

\textbf{【VHDL】}
\begin{lstlisting}
entity mux_logic is
  port (A, B, C : in std_logic; F : out std_logic);
end mux_logic;
architecture rtl of mux_logic is
  signal D : std_logic_vector(3 downto 0);
  signal S : std_logic_vector(1 downto 0);
begin
  S <= A & B;
  D <= '1' & (not C) & '0' & '1';  -- D3,D2,D1,D0
  with S select F <=
    D(0) when "00", D(1) when "01",
    D(2) when "10", D(3) when "11";
end rtl;
\end{lstlisting}

\begin{examfocus}
\textbf{这是组合逻辑的必考题！}两种考法：
(1) $2^n$选1 MUX实现n变量函数 → Dk直接接0/1
(2) $2^{n-1}$选1 MUX实现n变量函数 → 需要降维，剩余变量或其反接到D端
\end{examfocus}
\end{detailbox}


\section{变式2：双4选1 MUX构建8选1 MUX}

\begin{detailbox}[完整教学]
\textbf{【结构】}两片4选1 MUX处理低2位地址。一片2选1 MUX用最高位地址S2选择输出。
总MUX数：3片(2×4选1 + 1×2选1)。

\textbf{【VHDL——结构级】}
\begin{lstlisting}
-- 片0：S1S0选D0-D3
mux_low: entity work.mux4to1 port map(D=>D(3 downto 0), S=>S(1 downto 0), Y=>F_low);
-- 片1：S1S0选D4-D7
mux_high: entity work.mux4to1 port map(D=>D(7 downto 4), S=>S(1 downto 0), Y=>F_high);
-- 顶层2选1
F <= F_low when S(2)='0' else F_high;
\end{lstlisting}
\end{detailbox}


\section{变式3：带使能端的MUX}

\begin{detailbox}[完整教学]
\textbf{【功能】}E=1正常选择；E=0输出高阻态'Z'（三态输出）或固定'0'。

\textbf{【VHDL——三态输出】}
\begin{lstlisting}
F <= D(to_integer(unsigned(S))) when E='1' else 'Z';
\end{lstlisting}

\textbf{【应用】}多个MUX输出可共用一根总线（总线仲裁）。
\end{detailbox}


\section{变式4：MUX扩展——16选1}

\begin{detailbox}[完整教学]
\textbf{【树形结构】}4片4选1 MUX(第一层)+1片4选1 MUX(第二层)=16选1。
总MUX数：5片。

\textbf{【选择位分配】}S1S0→第一层(4片)，S3S2→第二层(1片)。
\end{detailbox}
